\chapter{Even a purely moral act (This Just Can't Be Right)}
One of my professors in college once gave an odd bit of advice: If you ever have to write on a work, and you're stuck for a topic, look for the exact midpoint, then write about whatever you find there. I am not remotely stuck on topics for Madoka Magica, but seeing as the end of this episode is the midpoint of the series, it seems as good a time as any to discuss Madoka Magica and consent issues.

The final scene of this episode has the characters recoiling in horror at the latest revelation from Kyubey: that the bodies of magical girls are not alive, but rather simply shells, which can be repaired so long as the Soul Gem is intact. Only by harming that gem can the magical girl herself be harmed—but by separating the gem from Sayaka's body, Madoka has effectively caused Sayaka's (temporary, thanks to Homura's quick intervention) death.

There is a case to be made (not a very good case, but a case nonetheless) that the girls are getting worked up over nothing. Frankly, what Kyubey describes seems like a pretty sweet deal: the physical experience of the body is close enough to being alive that most magical girls never even notice the change, but it is perfectly healthy, more durable, and able to heal from anything. Plus, given the evidence from Kyoko, it seems likely that it can eat junk food forever with no consequences. I'd take that deal in a heartbeat.

Indeed, the Soul Gems seem fairly clearly to be a reference to the Russian folkloric character Koshchyei Byessmyertnuy (Koschei the Deathless in English), who hid his soul in his finger, which he then severed and hid inside an egg inside a duck inside a hare inside an iron chest, which he buried underneath a green oak on a distant island. Koschei is a villainous figure who menaces young women, and only if the hero can find the egg can Koschei be harmed \cite{ref44}. The advantages to Koschei of doing this are quite clear, which is likely why his pop cultural heirs make use of similar protections (of which the Lich's phylactery in Dungeons \& Dragons \cite{ref45} and Voldemort's horcruxes in Harry Potter \cite{ref46} are likely the most widely known examples).

But there is an important difference between Koschei and Sayaka here, which is that of affirmative, informed consent. Koschei, the legends imply, knows what he is doing and chooses to do it. Sayaka had no idea that her life was in her Soul Gem, that her body had been transformed against her will. In a later episode she will note that she does not believe her new body is capable of bearing children, which she perhaps wanted to do someday. Regardless, the horror expressed by Sayaka, Kyoko, and Madoka in this episode makes it clear that all three recognize this as a supreme violation.

Kyubey's defense is that he doesn't understand why humans care so much about where their souls are located. This is irrelevant nonsense; it doesn't matter why they care when he clearly knows that they do care. Kyubey is deliberately concealing relevant information when he makes these pacts, and then blaming the victim when they reject him. In essence, he is justifying his actions by saying "Sayaka never said no."

"No means no" is often tossed around as a slogan in campaigns for women's rights, especially where issues of consent and bodily autonomy are concerned \cite{ref47}. However, while certainly better than failing to acknowledge that no means no, this is an incomplete standard, as Kyubey demonstrates. More important than "no means no" is "yes means yes," which is what is meant by a standard of affirmative consent \cite{ref48}. An absence of objection is insufficient, because that could mean that the person was unable to object, just as Sayaka was unable to object to aspects of the deal she didn't even know about.

This question of respecting the choices and autonomy of others interacts interestingly with another scene in this episode, when Madoka talks to her mother (in the vaguest possible terms, of course) about Sayaka's situation. Two things are important here, the first of which is Junko noting that doing the right thing does not always lead to happiness or good outcomes. The significance there is that it is an outright rejection of consequentialism (the meta-ethical view that the morality of an action is determined by its consequences \cite{ref49}) as an ethical position, which in general matches the stance taken by the show (hence the consistent depiction of Kyubey as a strong consequentialist).

The significance of rejecting consequentialism explicitly in the scene with Junko is that the ending scene on the bridge implicitly rejects it as well. Kyubey's position is a consequentialist one: the soul extraction is beneficial for the magical girls, since it enables them to fight witches and survive, but learning about it tends to make them unhappy, so the best thing to do is to extract the soul and not tell them about it. The music and the framing of the scene (particularly the way Kyubey is shot to be literally overshadowing the girls, despite his small stature) make it quite clear that the show is rejecting Kyubey's construction and empathizing with the girls' horror, which is to say rejecting the consequentialist perspective.

The second significant element of the conversation is the description of Sayaka as someone doing the right thing and making things worse as a consequence, because that description is hardly unique to Sayaka. It equally well applies to Homura, whose repeated attempts to save Madoka keep making her suffer more and become a more powerful witch in each successive timeline. Junko's advice to Madoka to make a mistake on her friend's behalf thus not only applies to throwing Sayaka's Soul Gem off the bridge; it is equally a description of her choice to become a magical girl (the very thing Homura has been trying to prevent) in the final episode. That this dual meaning is no accident seems clear given the musical choices; the music which accompanies Junko's advice in this episode, "Clementia," is reprised in the final episode's "Sagitta Luminis," which plays after Madoka makes her wish.

Further, this idea of saving someone by making a mistake for them is reiterated in Rebellion, where both Madoka and Homura take seemingly very ill-advised actions on each others' behalf—but more on that when we get there.
